\documentclass[11pt]{article}

\usepackage[a4paper,margin=1in]{geometry}
\usepackage{amsmath,amssymb,amsthm,mathtools}
\usepackage{enumitem}
\usepackage{hyperref}

\newtheorem{theorem}{Theorem}[section]
\newtheorem{lemma}[theorem]{Lemma}
\newtheorem{proposition}[theorem]{Proposition}
\newtheorem{definition}[theorem]{Definition}
\newtheorem{remark}[theorem]{Remark}

\DeclareMathOperator{\Tr}{Tr}
\DeclareMathOperator{\dist}{dist}
\newcommand{\R}{\mathbb{R}}
\newcommand{\bbP}{\mathbb{P}}
\newcommand{\bbE}{\mathbb{E}}
\newcommand{\OmegaGood}{\Omega}
\newcommand{\one}{\mathbf{1}}

\title{Appendix B: Plain-English Mathematical Translation of the Lean Proof}
\author{}
\date{}

\begin{document}
\maketitle

\section*{Purpose of this translation}

This note is a line-by-line mathematical translation of the Lean file
\[
\texttt{PptFactorization/AppendixB.lean}.
\]
It records exactly what the formal proof establishes.  The genuinely analytic
inputs of Appendix B are not reproved here; they remain explicit hypotheses,
as they do in Lean.  The verified content is the deterministic Lipschitz
bookkeeping, the good-set union bound, the comparison of the Levy exponent
with the desired scale, the median-to-mean reduction, the absorption of
numerical constants, and the final paper-facing wrappers.

\section{Local Lipschitz bookkeeping}

\begin{definition}[Local Lipschitz property]
Let \(X\) be a set, let \(\OmegaGood\subset X\), let
\(\rho:X\times X\to\R\) be a nonnegative distance-like function on
\(\OmegaGood\), and let \(f:X\to\R\).  We say that \(f\) is
\(L\)-Lipschitz on \(\OmegaGood\) if, for all \(x,y\in\OmegaGood\),
\[
  |f(x)-f(y)|\le L\,\rho(x,y).
\]
\end{definition}

\begin{lemma}[Enlarging a Lipschitz constant]
Assume \(\rho(x,y)\ge0\) for all \(x,y\in\OmegaGood\).  If \(f\) is
\(L\)-Lipschitz on \(\OmegaGood\) and \(L\le L'\), then \(f\) is
\(L'\)-Lipschitz on \(\OmegaGood\).
\end{lemma}

\begin{proof}
For \(x,y\in\OmegaGood\),
\[
  |f(x)-f(y)|\le L\rho(x,y)\le L'\rho(x,y),
\]
because \(\rho(x,y)\ge0\).  This is precisely the enlarged Lipschitz
estimate.
\end{proof}

\section{Deterministic matrix Lipschitz reductions}

The Lean proof first isolates the deterministic algebra which turns local
matrix estimates into a local Lipschitz estimate for a trace-power functional.
The matrix type is kept abstract: it only needs enough algebra to form powers
and differences.

\begin{proposition}[Trace perturbation plus difference bound]
Let \(A:X\to\mathcal A\) be a map into a ring, let
\(\tau:\mathcal A\to\R\) be the trace functional used in the paper, and let
\(\|\cdot\|_{1,*}:\mathcal A\to\R\) be an abstract trace-norm placeholder.
Fix \(k\in\mathbb N\).  Suppose that, for all \(x,y\in\OmegaGood\),
\[
 |\tau(A(x)^k)-\tau(A(y)^k)|
   \le a\,\|A(x)-A(y)\|_{1,*},
\]
and
\[
 \|A(x)-A(y)\|_{1,*}\le b\,\rho(x,y).
\]
Assume also \(a\ge0\), \(\rho(x,y)\ge0\) on \(\OmegaGood\), and
\[
  ab\le L.
\]
Then the function \(x\mapsto \tau(A(x)^k)\) is \(L\)-Lipschitz on
\(\OmegaGood\).
\end{proposition}

\begin{proof}
For \(x,y\in\OmegaGood\), multiply the difference bound by the nonnegative
number \(a\):
\[
 a\|A(x)-A(y)\|_{1,*}\le a\,b\,\rho(x,y).
\]
Hence
\[
\begin{aligned}
 |\tau(A(x)^k)-\tau(A(y)^k)|
 &\le a\|A(x)-A(y)\|_{1,*} \\
 &\le ab\,\rho(x,y) \\
 &\le L\,\rho(x,y),
\end{aligned}
\]
where the last step uses \(ab\le L\) and \(\rho(x,y)\ge0\).
\end{proof}

\begin{proposition}[Good-set size bound version]
Assume the preceding trace perturbation estimate holds with coefficient
\(a\).  Assume also that, on \(\OmegaGood\),
\[
  \|A(x)-A(y)\|_{1,*}
     \le D\,(\sigma(x)+\sigma(y))\,\rho(x,y),
\]
and that
\[
  \sigma(x)\le S \qquad (x\in\OmegaGood).
\]
If \(a\ge0\), \(D\ge0\), \(\rho\ge0\) on \(\OmegaGood\), and
\[
  a\,D\,(2S)\le L,
\]
then \(x\mapsto \tau(A(x)^k)\) is \(L\)-Lipschitz on \(\OmegaGood\).
\end{proposition}

\begin{proof}
For \(x,y\in\OmegaGood\), the size bound gives
\[
  \sigma(x)+\sigma(y)\le 2S.
\]
Multiplying by \(D\ge0\),
\[
  D(\sigma(x)+\sigma(y))\le D(2S).
\]
Thus
\[
 \|A(x)-A(y)\|_{1,*}
    \le D(2S)\rho(x,y).
\]
The previous proposition applies with \(b=D(2S)\).
\end{proof}

\begin{proposition}[Operator-norm version with trace norm]
Assume the following local matrix estimates on \(\OmegaGood\):
\[
 |\tau(A(x)^k)-\tau(A(y)^k)|
 \le k\,
   \max\{\|A(x)\|_{\mathrm{op}},\|A(y)\|_{\mathrm{op}}\}^{k-1}
   \|A(x)-A(y)\|_{1,*},
\]
\[
 \|A(x)\|_{\mathrm{op}}\le B,\qquad
 \|A(x)\|_{\mathrm{op}}\ge0,
\]
\[
 \|A(x)-A(y)\|_{1,*}
    \le D(\sigma(x)+\sigma(y))\rho(x,y),
\]
and
\[
 \sigma(x)\le S.
\]
Assume moreover that \(\|Z\|_{1,*}\ge0\) for all \(Z\), that \(B\ge0\), that
\(D\ge0\), and that
\[
   \bigl(k B^{k-1}\bigr)D(2S)\le L.
\]
Then \(x\mapsto \tau(A(x)^k)\) is \(L\)-Lipschitz on \(\OmegaGood\).
\end{proposition}

\begin{proof}
For \(x,y\in\OmegaGood\),
\[
 \max\{\|A(x)\|_{\mathrm{op}},\|A(y)\|_{\mathrm{op}}\}\le B.
\]
The maximum is nonnegative, so monotonicity of powers gives
\[
 \max\{\|A(x)\|_{\mathrm{op}},\|A(y)\|_{\mathrm{op}}\}^{k-1}
 \le B^{k-1}.
\]
Multiplying by \(k\ge0\), and then by the nonnegative difference norm, yields
\[
 k\max(\cdots)^{k-1}\|A(x)-A(y)\|_{1,*}
 \le
 kB^{k-1}\|A(x)-A(y)\|_{1,*}.
\]
Thus the trace perturbation estimate holds with coefficient \(a=kB^{k-1}\).
The good-set size version then gives the conclusion.
\end{proof}

\begin{proposition}[Frobenius/operator-norm version]
Let \(\|\cdot\|_{HS,*}\) be an abstract Frobenius or Hilbert--Schmidt norm.
Assume the local trace-power perturbation estimate
\[
 |\tau(A(x)^k)-\tau(A(y)^k)|
 \le
 \Delta\, k\,
 \max\{\|A(x)\|_{\mathrm{op}},\|A(y)\|_{\mathrm{op}}\}^{k-1}
 \|A(x)-A(y)\|_{HS,*},
\]
where \(\Delta\ge0\) is the deterministic dimension factor coming, in the
paper application, from an estimate of the form
\[
  |\Tr M|\le \Delta \|M\|_{HS}.
\]
Assume also
\[
 \|A(x)\|_{\mathrm{op}}\le B,\qquad
 \|A(x)\|_{\mathrm{op}}\ge0,
\]
\[
 \|A(x)-A(y)\|_{HS,*}
    \le D(\sigma(x)+\sigma(y))\rho(x,y),
\]
\[
 \sigma(x)\le S,
\]
with \(\|Z\|_{HS,*}\ge0\), \(B\ge0\), \(D\ge0\), and
\[
  \bigl(\Delta\,kB^{k-1}\bigr)D(2S)\le L.
\]
Then \(x\mapsto \tau(A(x)^k)\) is \(L\)-Lipschitz on \(\OmegaGood\).
\end{proposition}

\begin{proof}
The proof is identical to the preceding one, except that the trace coefficient
is
\[
  a=\Delta\,kB^{k-1}.
\]
The nonnegativity of \(\Delta\), \(k\), and \(B^{k-1}\) gives \(a\ge0\).
The operator-norm bound replaces
\[
 \Delta\,k\max(\cdots)^{k-1}
\]
by \(\Delta\,kB^{k-1}\).  The good-set size bound then converts the
Frobenius difference estimate into the Lipschitz bound
\[
 |\tau(A(x)^k)-\tau(A(y)^k)|\le L\rho(x,y).
\]
\end{proof}

\begin{remark}
This Frobenius version is the preferred deterministic route for the paper:
it avoids making the formalization depend on a full Schatten trace-norm API.
\end{remark}

\section{Good-set probability bookkeeping}

\begin{lemma}[Intersection good set]
Let \(\OmegaGood=\Omega_1\cap\Omega_2\).  Write
\[
 p_\Omega=\bbP(\OmegaGood^c),\qquad
 p_1=\bbP(\Omega_1^c),\qquad
 p_2=\bbP(\Omega_2^c).
\]
If
\[
 p_\Omega\le p_1+p_2,\qquad p_1\le E,\qquad p_2\le E,
\]
then
\[
 p_\Omega\le 2E.
\]
\end{lemma}

\begin{proof}
Immediate:
\[
  p_\Omega\le p_1+p_2\le E+E=2E.
\]
\end{proof}

\section{The Levy exponent scale}

\begin{lemma}[Exponential monotonicity]
If
\[
  c_2\,\frac{d^4\varepsilon^2}{k^2}
  \le
  \frac{n t^2}{4L^2},
\]
then
\[
  \exp\!\left(-\frac{n t^2}{4L^2}\right)
  \le
  \exp\!\left(-c_2\,\frac{d^4\varepsilon^2}{k^2}\right).
\]
\end{lemma}

\begin{proof}
The exponential function is increasing.  The assumed inequality becomes, after
multiplication by \(-1\),
\[
  -\frac{n t^2}{4L^2}
  \le
  -c_2\,\frac{d^4\varepsilon^2}{k^2}.
\]
Applying monotonicity of \(\exp\) gives the claim.
\end{proof}

\section{Median concentration}

\begin{proposition}[Localized Levy plus the good-set estimate]
Let \(M\) be a median reference value and let
\[
  \mathrm{Tail}_{\mathrm{med}}
\]
denote the corresponding median-centered deviation probability.  Suppose the
localized Levy lemma gives
\[
 \mathrm{Tail}_{\mathrm{med}}
 \le
 2\,\mathrm{bad}
 +2\exp\!\left(-\frac{nt^2}{4L^2}\right),
\]
and suppose the bad-set estimate gives
\[
 \mathrm{bad}\le 2\exp(-c d^2).
\]
If
\[
 c_2\,\frac{d^4\varepsilon^2}{k^2}
 \le
 \frac{nt^2}{4L^2},
\]
then
\[
 \mathrm{Tail}_{\mathrm{med}}
 \le
 4\exp(-c d^2)
 +2\exp\!\left(-c_2\,\frac{d^4\varepsilon^2}{k^2}\right).
\]
\end{proposition}

\begin{proof}
Insert the good-set bound into the localized Levy estimate:
\[
 2\,\mathrm{bad}\le 4\exp(-c d^2).
\]
The previous exponential comparison bounds the Levy exponential by the target
exponential.  Combining the two inequalities gives the displayed estimate.
\end{proof}

\section{Median-to-mean replacement}

\begin{proposition}[Mean--median gap from an integrated tail]
Let \((\Omega,\mu)\) be a probability space, let \(f:\Omega\to\R\) be
integrable, and let
\[
  m=\int_\Omega f\,d\mu.
\]
Fix a reference number \(M\).  Assume that
\[
  |f(\omega)-M|\le R
\]
for \(\mu\)-almost every \(\omega\), and that
\[
 \int_{(0,R]}
   \mu\{\omega:u\le |f(\omega)-M|\}\,du
 \le G.
\]
Then
\[
  |m-M|\le G.
\]
\end{proposition}

\begin{proof}
By Jensen's inequality in the form
\[
 \left|\int_\Omega (f-M)\,d\mu\right|
 \le
 \int_\Omega |f-M|\,d\mu,
\]
we have
\[
  |m-M|\le \int_\Omega |f-M|\,d\mu.
\]
Since \(0\le |f-M|\le R\) almost everywhere, the layer-cake identity gives
\[
 \int_\Omega |f-M|\,d\mu
 =
 \int_{(0,R]}
   \mu\{\omega:u\le |f(\omega)-M|\}\,du.
\]
The assumed integrated-tail bound therefore implies \(|m-M|\le G\).
\end{proof}

\begin{lemma}[Deterministic deviation inclusion]
Let \(z,m,M,s\in\R\).  If
\[
 |m-M|\le \frac{s}{2}
\]
and
\[
 s\le |z-m|,
\]
then
\[
 \frac{s}{2}\le |z-M|.
\]
\end{lemma}

\begin{proof}
The triangle inequality gives
\[
 |z-m|
 =
 |(z-M)+(M-m)|
 \le
 |z-M|+|M-m|.
\]
Since \(|M-m|=|m-M|\le s/2\), the inequality \(s\le |z-m|\) implies
\[
 s\le |z-M|+\frac{s}{2},
\]
hence \(s/2\le |z-M|\).
\end{proof}

\begin{lemma}[Probability version of the inclusion]
On any finite measure space, if \(|m-M|\le s/2\), then
\[
 \mu\{\omega:s\le |f(\omega)-m|\}
 \le
 \mu\{\omega:s/2\le |f(\omega)-M|\}.
\]
\end{lemma}

\begin{proof}
The deterministic lemma shows the first event is contained in the second
event.  Monotonicity of measure gives the probability inequality.
\end{proof}

\begin{proposition}[Median-to-mean from integrated centering]
Let \(m=\int f\,d\mu\), let \(f\) be integrable, and assume
\[
  |f-M|\le R
\]
almost everywhere.  If
\[
 \int_{(0,R]}
   \mu\{\omega:u\le |f(\omega)-M|\}\,du
 \le \frac{s}{2},
\]
then
\[
 \mu\{\omega:s\le |f(\omega)-m|\}
 \le
 \mu\{\omega:s/2\le |f(\omega)-M|\}.
\]
\end{proposition}

\begin{proof}
The integrated-tail estimate gives \(|m-M|\le s/2\) by the mean--median gap
proposition.  Apply the probability version of the deterministic inclusion.
\end{proof}

\begin{definition}[Quantitative median--mean bound]
The paper's displayed one-dimensional bound is encoded as
\[
  Q(R,\mathrm{bad},L,n)
  :=
  R\,\mathrm{bad}
  +\frac{L}{\sqrt n}.
\]
The term \(R\,\mathrm{bad}\) is the contribution of the bad set on the
deterministic range \(R\), and \(L/\sqrt n\) is the Gaussian integral scale
coming from the Levy tail.
\end{definition}

\begin{proposition}[Quantitative median-to-mean reduction]
If the integrated tail is bounded by
\[
 Q(R,\mathrm{bad},L,n)
\]
and
\[
 Q(R,\mathrm{bad},L,n)\le \frac{s}{2},
\]
then
\[
 \mu\{\omega:s\le |f(\omega)-m|\}
 \le
 \mu\{\omega:s/2\le |f(\omega)-M|\}.
\]
\end{proposition}

\begin{proof}
The two inequalities imply that the integrated tail is at most \(s/2\), so the
previous proposition applies.
\end{proof}

\section{Mean-centered concentration}

\begin{proposition}[Mean concentration from a centering input]
Assume
\[
 \mathrm{Tail}_{\mathrm{mean}}
 \le
 \mathrm{Tail}_{\mathrm{med}},
\]
and assume the localized Levy estimate, the good-set bound, and the exponent
comparison from the median concentration proposition.  Then
\[
 \mathrm{Tail}_{\mathrm{mean}}
 \le
 4\exp(-c d^2)
 +2\exp\!\left(-c_2\,\frac{d^4\varepsilon^2}{k^2}\right).
\]
\end{proposition}

\begin{proof}
The median concentration proposition gives the displayed upper bound for
\(\mathrm{Tail}_{\mathrm{med}}\).  Since
\(\mathrm{Tail}_{\mathrm{mean}}\le\mathrm{Tail}_{\mathrm{med}}\), the same
bound holds for the mean-centered tail.
\end{proof}

\begin{proposition}[Mean concentration from integrated centering]
Let \(m=\int f\,d\mu\).  Suppose the integrated median tail is at most
\(s/2\), so that
\[
 \mu\{\omega:s\le |f(\omega)-m|\}
 \le
 \mu\{\omega:s/2\le |f(\omega)-M|\}.
\]
Suppose also that the localized Levy lemma bounds the half-scale median tail:
\[
 \mu\{\omega:s/2\le |f(\omega)-M|\}
 \le
 2\,\mathrm{bad}
 +2\exp\!\left(-\frac{nt^2}{4L^2}\right).
\]
If
\[
 \mathrm{bad}\le 2\exp(-c d^2)
\]
and
\[
 c_2\,\frac{d^4\varepsilon^2}{k^2}
 \le
 \frac{nt^2}{4L^2},
\]
then
\[
 \mu\{\omega:s\le |f(\omega)-m|\}
 \le
 4\exp(-c d^2)
 +2\exp\!\left(-c_2\,\frac{d^4\varepsilon^2}{k^2}\right).
\]
\end{proposition}

\begin{proof}
The integrated centering input gives the inclusion of the mean-centered event
inside the half-scale median-centered event.  Therefore the mean-centered
tail is bounded by the median tail appearing in the localized Levy estimate.
The good-set and exponent comparisons then give the displayed estimate,
exactly as in the median concentration proposition.
\end{proof}

\begin{lemma}[Absorbing numerical constants]
Suppose
\[
  T\le 4E_1+2E_2,
\]
where \(E_1,E_2\ge0\).  If \(K\ge4\), then
\[
  T\le K E_1+K E_2.
\]
\end{lemma}

\begin{proof}
Since \(K\ge4\), also \(K\ge2\).  Hence
\[
  4E_1\le KE_1,\qquad 2E_2\le KE_2,
\]
and summing gives the claim.
\end{proof}

\begin{theorem}[Final abstract concentration theorem]
Under the centering input, the localized Levy estimate, the good-set estimate,
the exponent comparison, and \(K\ge4\),
\[
 \mathrm{Tail}_{\mathrm{mean}}
 \le
 K\exp(-c d^2)
 +K\exp\!\left(-c_2\,\frac{d^4\varepsilon^2}{k^2}\right).
\]
\end{theorem}

\begin{proof}
First prove the explicit estimate with coefficients \(4\) and \(2\).  Then
use the absorption lemma with
\[
 E_1=\exp(-c d^2),
 \qquad
 E_2=\exp\!\left(-c_2\,\frac{d^4\varepsilon^2}{k^2}\right).
\]
Both exponentials are nonnegative.
\end{proof}

\section{Paper-facing notation}

The Lean file then introduces notation matching the manuscript.

\begin{definition}[Sphere dimension]
\[
  N(d,s):=2d^2s.
\]
This is the real dimension of \(M_{d^2,s}(\mathbb C)\) viewed as a real
Hilbert space.
\end{definition}

\begin{definition}[Natural deviation scale]
For a moment order \(r\in\mathbb N\),
\[
  a_{d,\varepsilon,r}
  :=
  \frac{\varepsilon}{d^{2r-2}}.
\]
\end{definition}

\begin{definition}[Local Lipschitz scale]
For a constant \(C\) and moment parameter \(k\),
\[
  L_{C,k,d,r}
  :=
  \frac{Ck}{d^{2r-2}}.
\]
\end{definition}

\begin{definition}[Levy and target exponents]
\[
  \mathcal L(n,t,L):=\frac{nt^2}{4L^2},
\]
and
\[
  \mathcal T(c_{\mathrm{loc}},d,\varepsilon,k)
  :=
  c_{\mathrm{loc}}\frac{d^4\varepsilon^2}{k^2}.
\]
\end{definition}

\begin{definition}[Good-set and final paper tails]
\[
  G(c_{\mathrm{sph}},d):=\exp(-c_{\mathrm{sph}}d^2),
\]
\[
  B_{\mathrm{med}}
  :=
  4G(c_{\mathrm{sph}},d)
  +2\exp(-\mathcal T(c_{\mathrm{loc}},d,\varepsilon,k)),
\]
and
\[
  B_{\mathrm{paper}}
  :=
  K G(c_{\mathrm{sph}},d)
  +K\exp(-\mathcal T(c_{\mathrm{loc}},d,\varepsilon,k)).
\]
\end{definition}

\begin{lemma}[Paper notation for the event inclusion]
If
\[
 |m-M|\le \frac{a_{d,\varepsilon,r}}{2}
\]
and
\[
 a_{d,\varepsilon,r}\le |z-m|,
\]
then
\[
 \frac{a_{d,\varepsilon,r}}{2}\le |z-M|.
\]
\end{lemma}

\begin{proof}
This is the deterministic deviation inclusion with
\(s=a_{d,\varepsilon,r}\).
\end{proof}

\begin{lemma}[Paper notation for the good set]
If
\[
 \mathrm{bad}\le \mathrm{bad}_1+\mathrm{bad}_2,\qquad
 \mathrm{bad}_1\le G(c_{\mathrm{sph}},d),\qquad
 \mathrm{bad}_2\le G(c_{\mathrm{sph}},d),
\]
then
\[
 \mathrm{bad}\le 2G(c_{\mathrm{sph}},d).
\]
\end{lemma}

\begin{proof}
This is the good-set intersection bound.
\end{proof}

\begin{proposition}[Paper notation for median concentration]
If
\[
 \mathrm{Tail}_{\mathrm{med}}
 \le
 2\,\mathrm{bad}
 +2\exp(-\mathcal L(n,t,L)),
\]
\[
 \mathrm{bad}\le 2G(c_{\mathrm{sph}},d),
\]
and
\[
 \mathcal T(c_{\mathrm{loc}},d,\varepsilon,k)\le \mathcal L(n,t,L),
\]
then
\[
 \mathrm{Tail}_{\mathrm{med}}\le B_{\mathrm{med}}.
\]
\end{proposition}

\begin{proof}
Unfolding the definitions gives exactly the median concentration proposition.
\end{proof}

\section{The scale computations}

\begin{proposition}[Exact-scale exponent comparison]
Assume
\[
  d\ne0,\qquad C\ne0,\qquad k\ne0,
\]
and
\[
  c_{\mathrm{loc}}\le \frac{c_{\mathrm{dim}}}{4C^2}.
\]
Then
\[
 \mathcal T(c_{\mathrm{loc}},d,\varepsilon,k)
 \le
 \mathcal L\!\left(
     c_{\mathrm{dim}}d^4,\,
     \frac{\varepsilon}{d^{2r-2}},\,
     \frac{Ck}{d^{2r-2}}
   \right).
\]
\end{proposition}

\begin{proof}
Set
\[
  q=d^{2r-2}.
\]
Since \(d\ne0\), \(q\ne0\).  Direct algebra gives
\[
\begin{aligned}
 \mathcal L\!\left(
     c_{\mathrm{dim}}d^4,\,
     \frac{\varepsilon}{q},\,
     \frac{Ck}{q}
   \right)
 &=
 \frac{(c_{\mathrm{dim}}d^4)(\varepsilon/q)^2}
      {4(Ck/q)^2} \\
 &=
 \frac{c_{\mathrm{dim}}}{4C^2}
 \frac{d^4\varepsilon^2}{k^2}.
\end{aligned}
\]
The factor \(d^4\varepsilon^2/k^2\) is nonnegative.  Multiplying the
assumption
\[
 c_{\mathrm{loc}}\le \frac{c_{\mathrm{dim}}}{4C^2}
\]
by this nonnegative factor gives
\[
 c_{\mathrm{loc}}\frac{d^4\varepsilon^2}{k^2}
 \le
 \frac{c_{\mathrm{dim}}}{4C^2}\frac{d^4\varepsilon^2}{k^2},
\]
which is the desired comparison.
\end{proof}

\begin{proposition}[Half-scale exponent comparison]
Assume
\[
  d\ne0,\qquad C\ne0,\qquad k\ne0,
\]
and
\[
  c_{\mathrm{loc}}\le \frac{c_{\mathrm{dim}}}{16C^2}.
\]
Then
\[
 \mathcal T(c_{\mathrm{loc}},d,\varepsilon,k)
 \le
 \mathcal L\!\left(
     c_{\mathrm{dim}}d^4,\,
     \frac{1}{2}\frac{\varepsilon}{d^{2r-2}},\,
     \frac{Ck}{d^{2r-2}}
   \right).
\]
\end{proposition}

\begin{proof}
Again put \(q=d^{2r-2}\).  Then
\[
\begin{aligned}
 \mathcal L\!\left(
     c_{\mathrm{dim}}d^4,\,
     \frac{\varepsilon/q}{2},\,
     \frac{Ck}{q}
   \right)
 &=
 \frac{(c_{\mathrm{dim}}d^4)((\varepsilon/q)/2)^2}
      {4(Ck/q)^2} \\
 &=
 \frac{c_{\mathrm{dim}}}{16C^2}
 \frac{d^4\varepsilon^2}{k^2}.
\end{aligned}
\]
Multiplying
\[
 c_{\mathrm{loc}}\le \frac{c_{\mathrm{dim}}}{16C^2}
\]
by the nonnegative factor \(d^4\varepsilon^2/k^2\) proves the claim.
\end{proof}

\section{Paper-facing final wrappers}

\begin{theorem}[Appendix B with abstract paper inputs]
Suppose the following data are given:
\[
 \mathrm{Tail}_{\mathrm{mean}},\quad
 \mathrm{Tail}_{\mathrm{med}},\quad
 \mathrm{bad},\quad
 n,\quad t,\quad L,\quad
 c_{\mathrm{sph}},\quad c_{\mathrm{loc}},\quad d,\quad
 \varepsilon,\quad k,\quad K.
\]
Assume
\[
 \mathrm{Tail}_{\mathrm{mean}}\le \mathrm{Tail}_{\mathrm{med}},
\]
\[
 \mathrm{Tail}_{\mathrm{med}}
 \le
 2\,\mathrm{bad}
 +2\exp(-\mathcal L(n,t,L)),
\]
\[
 \mathrm{bad}\le 2G(c_{\mathrm{sph}},d),
\]
\[
 \mathcal T(c_{\mathrm{loc}},d,\varepsilon,k)\le \mathcal L(n,t,L),
\]
and \(K\ge4\).  Then
\[
 \mathrm{Tail}_{\mathrm{mean}}
 \le
 B_{\mathrm{paper}}.
\]
\end{theorem}

\begin{proof}
Unfolding \(B_{\mathrm{paper}}\), \(G\), \(\mathcal T\), and \(\mathcal L\),
this is exactly the final abstract concentration theorem.
\end{proof}

\begin{theorem}[Two-branch final wrapper]
Assume \(\mathrm{Tail}_{\mathrm{mean}}\le1\).  Suppose that either
\[
 \mathrm{Tail}_{\mathrm{mean}}\le \mathrm{Tail}_{\mathrm{med}},
\]
or
\[
 1\le B_{\mathrm{paper}}.
\]
Assume also the localized Levy estimate, the good-set estimate, the exponent
comparison, and \(K\ge4\).  Then
\[
  \mathrm{Tail}_{\mathrm{mean}}\le B_{\mathrm{paper}}.
\]
\end{theorem}

\begin{proof}
There are two cases.

In the first case,
\(\mathrm{Tail}_{\mathrm{mean}}\le \mathrm{Tail}_{\mathrm{med}}\), the
abstract paper-input theorem applies directly.

In the second case, \(1\le B_{\mathrm{paper}}\).  Since
\(\mathrm{Tail}_{\mathrm{mean}}\le1\), transitivity gives
\[
  \mathrm{Tail}_{\mathrm{mean}}\le1\le B_{\mathrm{paper}}.
\]
\end{proof}

\begin{theorem}[Wrapper with displayed paper scales]
Assume the mean-to-median centering input, the localized Levy estimate with
\[
 n=N(d,s),\qquad
 t=a_{d,\varepsilon,r},\qquad
 L=L_{C,k,d,r},
\]
the good-set estimate, the corresponding exponent comparison, and \(K\ge4\).
Then
\[
 \mathrm{Tail}_{\mathrm{mean}}
 \le
 B_{\mathrm{paper}}.
\]
\end{theorem}

\begin{proof}
This is the abstract paper-input theorem with
\[
 n=N(d,s),\qquad
 t=a_{d,\varepsilon,r},\qquad
 L=L_{C,k,d,r}.
\]
\end{proof}

\begin{theorem}[Wrapper from two good sets and exact scales]
Assume
\[
 d\ne0,\qquad C\ne0,\qquad k\ne0,\qquad
 c_{\mathrm{loc}}\le \frac{c_{\mathrm{dim}}}{4C^2}.
\]
Assume the centering input
\[
 \mathrm{Tail}_{\mathrm{mean}}\le \mathrm{Tail}_{\mathrm{med}},
\]
the two good-set bounds
\[
 \mathrm{bad}\le \mathrm{bad}_1+\mathrm{bad}_2,\qquad
 \mathrm{bad}_1\le G(c_{\mathrm{sph}},d),\qquad
 \mathrm{bad}_2\le G(c_{\mathrm{sph}},d),
\]
the localized Levy estimate with
\[
 n=c_{\mathrm{dim}}d^4,\qquad
 t=a_{d,\varepsilon,r},\qquad
 L=L_{C,k,d,r},
\]
and \(K\ge4\).  Then
\[
 \mathrm{Tail}_{\mathrm{mean}}\le B_{\mathrm{paper}}.
\]
\end{theorem}

\begin{proof}
The two one-set estimates and the union bound give
\[
  \mathrm{bad}\le 2G(c_{\mathrm{sph}},d).
\]
The exact-scale computation proves
\[
 \mathcal T(c_{\mathrm{loc}},d,\varepsilon,k)
 \le
 \mathcal L(c_{\mathrm{dim}}d^4,a_{d,\varepsilon,r},L_{C,k,d,r}).
\]
The abstract paper-input theorem then applies.
\end{proof}

\begin{theorem}[Strongest integrated-centering wrapper formalized in Lean]
Let \((\Omega,\mu)\) be a probability space, let \(f:\Omega\to\R\), and put
\[
 m=\int_\Omega f\,d\mu.
\]
Fix a median reference value \(M\).  Assume
\[
 d\ne0,\qquad C\ne0,\qquad k\ne0,\qquad
 c_{\mathrm{loc}}\le \frac{c_{\mathrm{dim}}}{16C^2}.
\]
Assume \(f\) is integrable, \(|f-M|\le R\) almost everywhere, and
\[
 \int_{(0,R]}
   \mu\{\omega:u\le |f(\omega)-M|\}\,du
 \le
 \frac{a_{d,\varepsilon,r}}{2}.
\]
Assume the two good-set estimates
\[
 \mathrm{bad}\le \mathrm{bad}_1+\mathrm{bad}_2,\qquad
 \mathrm{bad}_1\le G(c_{\mathrm{sph}},d),\qquad
 \mathrm{bad}_2\le G(c_{\mathrm{sph}},d),
\]
and the localized Levy estimate at half scale:
\[
 \mu\left\{\omega:
   \frac{a_{d,\varepsilon,r}}{2}\le |f(\omega)-M|
 \right\}
 \le
 2\,\mathrm{bad}
 +2\exp\!\left(
   -\mathcal L\left(
       c_{\mathrm{dim}}d^4,\,
       \frac{a_{d,\varepsilon,r}}{2},\,
       L_{C,k,d,r}
     \right)
 \right).
\]
If \(K\ge4\), then
\[
 \mu\{\omega:a_{d,\varepsilon,r}\le |f(\omega)-m|\}
 \le
 B_{\mathrm{paper}}.
\]
\end{theorem}

\begin{proof}
First, the two good-set estimates imply
\[
  \mathrm{bad}\le 2G(c_{\mathrm{sph}},d).
\]
Second, the half-scale exponent computation gives
\[
 \mathcal T(c_{\mathrm{loc}},d,\varepsilon,k)
 \le
 \mathcal L\left(
       c_{\mathrm{dim}}d^4,\,
       \frac{a_{d,\varepsilon,r}}{2},\,
       L_{C,k,d,r}
     \right).
\]
Third, the integrated-tail assumption gives
\[
 |m-M|\le \frac{a_{d,\varepsilon,r}}{2}.
\]
Therefore
\[
 \left\{\omega:a_{d,\varepsilon,r}\le |f(\omega)-m|\right\}
 \subset
 \left\{\omega:\frac{a_{d,\varepsilon,r}}{2}\le |f(\omega)-M|\right\}.
\]
The half-scale localized Levy estimate, the good-set bound, and the exponent
comparison give the explicit estimate
\[
 \mu\{\omega:a_{d,\varepsilon,r}\le |f(\omega)-m|\}
 \le
 4G(c_{\mathrm{sph}},d)
 +2\exp(-\mathcal T(c_{\mathrm{loc}},d,\varepsilon,k)).
\]
Finally, since \(K\ge4\) and both exponential terms are nonnegative, the
absorption lemma converts this into
\[
 \mu\{\omega:a_{d,\varepsilon,r}\le |f(\omega)-m|\}
 \le
 K G(c_{\mathrm{sph}},d)
 +K\exp(-\mathcal T(c_{\mathrm{loc}},d,\varepsilon,k)),
\]
which is \(B_{\mathrm{paper}}\).
\end{proof}

\section*{What remains analytic}

The Lean proof, and therefore this translation, does not prove the following
analytic ingredients.  They are supplied as hypotheses:
\begin{enumerate}[label=\arabic*.]
  \item The trace-power perturbation inequality and its concrete matrix-norm
  realization.
  \item The operator/Frobenius norm estimates defining the good set.
  \item The two one-set probability bounds for \(\Omega_1^c\) and
  \(\Omega_2^c\).
  \item The localized Levy lemma on the Hilbert--Schmidt sphere.
  \item The one-dimensional integrated-tail estimate used to pass from the
  median to the mean.
\end{enumerate}

All remaining steps are deterministic or probability-bookkeeping steps, and
are exactly the steps certified by the Lean file.

\end{document}
